
\documentclass{article}
\usepackage{colortbl}
\usepackage{makecell}
\usepackage{multirow}
\usepackage{supertabular}

\begin{document}

\newcounter{utterance}

\centering \large Interaction Transcript for game `cladder', experiment `full\_v1.5\_default', episode 2017 with BSU{-}SLIM\_\_Qwen3.5{-}9B{-}Base\_\_playpen{-}prm{-}9b{-}best\_of\_n.
\vspace{24pt}

{ \footnotesize  \setcounter{utterance}{1}
\setlength{\tabcolsep}{0pt}
\begin{supertabular}{c@{$\;$}|p{.15\linewidth}@{}p{.15\linewidth}p{.15\linewidth}p{.15\linewidth}p{.15\linewidth}p{.15\linewidth}}
    \# & \multicolumn{2}{c}{Player} && \multicolumn{2}{c}{Game Master} \\
    \hline

    \theutterance \stepcounter{utterance}  
    & & & \multicolumn{4}{p{0.6\linewidth}}{
        \cellcolor[rgb]{0.9,0.9,0.9}{
            \makecell[{{p{\linewidth}}}]{
                \texttt{\tiny{[P1$\langle$GM]}}
                \texttt{You are an expert in causal inference. The following question is not a typical commonsense query, but rather a meticulously designed question created by a professor specializing in causal inference, intended to assess the students' mastery of the course content.} \\
\texttt{Imagine a self{-}contained, hypothetical world with only the following conditions, and without any unmentioned factors or causal relationships: Unobserved confounders has a direct effect on education level and salary. Proximity to a college has a direct effect on education level. Education level has a direct effect on salary. Unobserved confounders is unobserved. For people living far from a college, the probability of high salary is 39\%. For people living close to a college, the probability of high salary is 49\%. For people living far from a college, the probability of college degree or higher is 57\%. For people living close to a college, the probability of college degree or higher is 30\%. Will college degree or higher increase the chance of high salary?} \\
\texttt{Start your answer with 'Yes' or 'No', followed by additional reasoning or evidence to support your explanation.} \\
            }
        }
    }
     \\ \\

    \theutterance \stepcounter{utterance}  
    & \multicolumn{4}{p{0.6\linewidth}}{
        \cellcolor[rgb]{0.561,0.737,0.561}{
            \makecell[{{p{\linewidth}}}]{
                \texttt{\tiny{[P1$\rangle$GM]}}
                \texttt{No} \\
\\ 
\texttt{**Reasoning:**} \\
\\ 
\texttt{To determine causality in this world, we must analyze whether the observed association between "college degree" and "high salary" is confounded. In this hypothetical scenario, "proximity to a college" acts as a proxy for "unobserved confounders," but crucially, proximity also has a **direct effect** on education level.} \\
\\ 
\texttt{Let's break down the causal structure and the direction of the effects:} \\
\texttt{1.  **Causal Pathways:**} \\
\texttt{    *   Unobserved Confounders (\$U\$) \$\textbackslash{}rightarrow\$ Education (\$E\$)} \\
\texttt{    *   Unobserved Confounders (\$U\$) \$\textbackslash{}rightarrow\$ Salary (\$S\$)} \\
\texttt{    *   Proximity to College (\$P\$) \$\textbackslash{}rightarrow\$ Education (\$E\$)} \\
\texttt{    *   Education (\$E\$) \$\textbackslash{}rightarrow\$ Salary (\$S\$)} \\
\\ 
\texttt{2.  **Observed Data (Intervention on Proximity):**} \\
\texttt{    *   **Far from college (\$P=0\$):** Probability of College Degree (\$E=1\$) = 57\%. Probability of High Salary (\$S=1\$) = 39\%.} \\
\texttt{    *   **Close to college (\$P=1\$):** Probability of College Degree (\$E=1\$) = 30\%. Probability of High Salary (\$S=1\$) = 49\%.} \\
\\ 
\texttt{3.  **Analysis of the Correlation vs. Causation:**} \\
\texttt{    *   **Observed Association:** Comparing the two groups, those living **close** to a college have a **lower** probability of getting a degree (30\% vs 57\%) but a **higher** probability of high salary (49\% vs 39\%). The raw correlation shows that getting a degree is associated with a *lower* salary in this specific dataset.} \\
\\ 
\texttt{4.  **Decomposing the Bias (Backdoor Path):**} \\
\texttt{    *   The bias arises because Proximity (\$P\$) affects Education (\$E\$).} \\
\texttt{    *   In this hypothetical world, the mechanism of \$P\$ on \$E\$ is inverse: Living close to a college reduces the likelihood of obtaining a degree (perhaps due to lack of competition, or a different population selecting into those areas, or simply the defined direct effect in this specific problem). Moving from Far (\$E=57\textbackslash{}\%\$) to Close (\$E=30\textbackslash{}\%\$) causes a **decrease** in the probability of having a degree.} \\
\texttt{    *   Simultaneously, look at the effect of \$P\$ on Salary (\$S\$) *independent* of Education. Since \$U\$ affects both \$E\$ and \$S\$, and \$P\$ is a proxy for \$U\$ (or correlated with it), we must look at the marginal effect.} \\
\texttt{    *   However, the question asks: "Will {[}getting{]} college degree or higher increase the chance of high salary?" This asks for the **causal effect** of \$E\$, not the correlation between \$P\$ and \$S\$.} \\
\texttt{    *   To find the causal effect of \$E\$, we must block the backdoor path \$P \textbackslash{}rightarrow E \textbackslash{}rightarrow S\$ by comparing groups with the *same* \$P\$ but different \$E\$, OR we must mathematically adjust for the confounding factor \$P\$.} \\
\texttt{    *   Since \$P\$ is unobserved, we cannot adjust for it using standard regression. We must use the experimental data provided (which effectively compares the two levels of \$P\$) to infer the direction of the causal arrows or recognize the specific bias introduced.} \\
\\ 
\texttt{5.  **Logical Deduction in this Specific World:**} \\
\texttt{    *   Let's assume a simple linear probability model or monotonicity to test the logic.} \\
\texttt{    *   Let \$P(E|P=Far) = 0.57\$ and \$P(S|P=Far) = 0.39\$.} \\
\texttt{    *   Let \$P(E|P=Close) = 0.30\$ and \$P(S|P=Close) = 0.49\$.} \\
\texttt{    *   Notice the trade{-}off: High proximity reduces the chance of a degree but increases the chance of high salary.} \\
\texttt{    *   Why? Because Proximity (\$P\$) has a direct effect on Salary *independent* of Education (since \$P \textbackslash{}rightarrow U \textbackslash{}rightarrow S\$ exists, and \$P \textbackslash{}rightarrow E\$ exists). Wait, the prompt says "Unobserved confounders has a direct effect on education level and salary." It does *not* explicitly state that Proximity has a direct effect on salary, only that it has a direct effect on education.} \\
\texttt{    *   Let's re{-}read carefully: "Proximity to a college has a direct effect on education level." It does not say Proximity affects salary directly.} \\
\texttt{    *   So the graph is: \$U \textbackslash{}rightarrow} \\
            }
        }
    }
    & & \\ \\

    \theutterance \stepcounter{utterance}  
    & & & \multicolumn{2}{p{0.3\linewidth}}{
        \cellcolor[rgb]{0.9,0.9,0.9}{
            \makecell[{{p{\linewidth}}}]{
                \texttt{\tiny{[GM$|$GM]}}
                \texttt{No} \\
            }
        }
    }
    & & \\ \\

    \theutterance \stepcounter{utterance}  
    & & & \multicolumn{2}{p{0.3\linewidth}}{
        \cellcolor[rgb]{0.9,0.9,0.9}{
            \makecell[{{p{\linewidth}}}]{
                \texttt{\tiny{[GM$|$GM]}}
                \texttt{no} \\
            }
        }
    }
    & & \\ \\

    \theutterance \stepcounter{utterance}  
    & & & \multicolumn{2}{p{0.3\linewidth}}{
        \cellcolor[rgb]{0.9,0.9,0.9}{
            \makecell[{{p{\linewidth}}}]{
                \texttt{\tiny{[GM$|$GM]}}
                \texttt{game\_result = WIN} \\
            }
        }
    }
    & & \\ \\

\end{supertabular}
}

\end{document}
